% Este arquivo contém o conteúdo da seção de Revisão Sistemática da Literatura (Systematic Review).
% Ele será incluído no arquivo principal usando \secao{Systematic Review}{partes/systematic_review}.

This section outlines the methodology employed for the Systematic Literature Review (SLR) conducted to identify, evaluate, and synthesize pertinent studies concerning the parallelization and optimization of the Aho-Corasick algorithm. The review focuses on implementations targeting multi-core CPUs. The objective of this SLR is to support the development of the present research and to answer the formulated research questions.

\subsecao{Research Questions}

The following Research Questions (RQs) were formulated to guide the literature search, selection, and analysis process:
\begin{enumerate}
    \item \textbf{RQ1:} What parallelization techniques and associated challenges are documented in the literature for implementing the Aho-Corasick algorithm utilizing threads on multi-core CPUs?
    \item \textbf{RQ2:} What specific optimizations for the Aho-Corasick algorithm are employed in application scenarios involving extensive pattern sets, particularly within domains of information security?
\end{enumerate}

\subsecao{Search Strategy and Study Selection Protocol}

The search and selection protocol for identifying primary studies encompassed the definition of data sources, construction of search queries, and establishment of inclusion and exclusion criteria.

\subsubsecao{Data Sources}
The following scientific databases were selected due to their scope and relevance in the area of ​​Computer Science and Software Engineering:
\begin{itemize}
    \item Web of Science (WoS)
    \item Scopus
    \item ACM Digital Library (ACM DL)
    \item IEEE Xplore
\end{itemize}

\subsubsecao{Search Queries}

Two primary search strings were constructed to address the aspects of parallelization (RQ1) and optimization/application (RQ2) pertinent to the Aho-Corasick algorithm. The literature search was executed in May 2025.

\begin{tabular}{@{} p{\linewidth-2em} @{}} % Coluna com largura da linha menos um pequeno recuo

    \textbf{String 1 (Focus on Parallelization):} \\
    \texttt{"Aho-Corasick" AND (parallel* OR multithread* OR multi-thread* OR thread* OR concurren* OR hpc OR high-perform* OR high perform*) AND ("multi-core" OR multicore OR "many-core" OR CPU) NOT "distributed systems")} \\

    \textbf{String 2 (Focus on Optimization and Applications):} \\
    \texttt{"Aho-Corasick" AND (optimi* OR accelerat* OR perform* OR evaluat* OR improve*) AND (“application” OR “implementation” OR “use case”))} \\

\end{tabular}

\subsubsecao{Inclusion (IC) and Exclusion (EC) Criteria}

The subsequent criteria were established for the selection of primary studies:

\textbf{Inclusion Criteria (IC):}
\begin{itemize}
    \item \textbf{IC1 (Primary Focus):} The study addresses the Aho-Corasick (AC) algorithm or presents parallelization/optimization techniques demonstrably applicable to AC.
    \item \textbf{IC2 (Parallelism Aspect):} The study investigates, proposes, implements, analyzes, or explicitly discusses the parallelization or concurrent execution of the Aho-Corasick algorithm.
    \item \textbf{IC3 (Target Platform):} The study reports empirical results or analyses about multi-core CPU platforms.
    \item \textbf{IC4 (RQ Relevance):} The study provides information relevant to addressing at least one RQ.
    \item \textbf{IC5 (Language):} Published in English or Portuguese.
    \item \textbf{IC6 (Accessibility):} Full-text availability of the study.
    \item \textbf{IC7 (Publication Window):} Published between January 2015 and May 2025, inclusive.
\end{itemize}

\textbf{Exclusion Criteria (EC):}
\begin{itemize}
    \item \textbf{EC1 (Irrelevant Focus):} Studies with only superficial mentions of the Aho-Corasick algorithm or pattern matching, lacking a primary focus on its parallelization or optimization.
    \item \textbf{EC2 (Non-CPU Platform Exclusivity):} Studies exclusively focused on non-CPU platforms (e.g., GPU, FPGA) without analysis or insights transferable to multi-core CPU environments.
    \item \textbf{EC3 (Insufficient Detail):} Studies lacking sufficient detail regarding methodology, implementation, or reported results.
    \item \textbf{EC4 (Publication Type):} Opinion pieces, editorials, and non-peer-reviewed gray literature.
    \item \textbf{EC5 (Language Barrier):} Publications in languages other than English or Portuguese.
    \item \textbf{EC6 (Duplication):} Duplicate publications (the most comprehensive or recent version was retained).
    \item \textbf{EC7 (Inaccessibility):} Full-text unobtainable.
\end{itemize}

\subsubsecao{Selection Process and Initial Search Outcomes}

The study selection was a multi-stage process. Initially, the constructed search strings were executed against the selected digital libraries. The retrieved results were subsequently refined using filters provided by each database (e.g., publication period, document type, language, subject area), as detailed in Table \ref{tab:search_results}:

\begin{table}[ht]
    \centering
    \caption{Search results and filters applied per database.}
    \label{tab:search_results}
    \small
    \begin{tabular}{@{} l c c p{7.5cm} @{}}
        \toprule
        \textbf{Database} & \textbf{String 1} & \textbf{String 2} & \textbf{Filters Applied} \\
        & \multicolumn{1}{c}{\scriptsize{(Found / Kept)}} & \multicolumn{1}{c}{\scriptsize{(Found / Kept)}} & \\
        \midrule
        ACM Digital Library & 122 / 29 & 274 / 66 & 2015-2025; Research Article; Publisher: ACM; Title excl.: GPU; Lang: English. \\ \addlinespace
        Web of Science & 30 / 6 & 66 / 20 & 2015-2025; WoS Cats: CS (excl. AI); Title excl.: GPU; Lang: English. \\ \addlinespace
        Scopus & 34 / 6 & 145 / 49 & 2015-2025; Lang: English; Subj: CS, Eng; Types: Article, Conf. Paper; Title excl.: GPU. \\ \addlinespace
        IEEE Xplore & 25 / 2 & 60 / 18 & 2015-2025; Title excl.: GPU; Lang: English. \\
        \bottomrule
    \end{tabular}
\end{table}

The \texttt{Title excl.: GPU} filter listed in Table~\ref{tab:search_results} acted only as a coarse identification-stage noise reducer, not as the formal platform criterion: the platform decision is governed by EC2, under which studies on non-CPU platforms (GPU, FPGA, ASIC, SIMD) were retained whenever they offered techniques or findings transferable to shared-memory multi-core CPUs, as argued individually in each review below.

Following the initial database search and filtering stage, the resultant sets of candidate articles were aggregated. The 196 records retained across both search strings and all four databases (Table~\ref{tab:search_results}) were merged, and duplicate entries were removed, yielding 24 unique studies. Each of these was screened by reading its abstract and conclusion against the inclusion and exclusion criteria; the 10 studies most directly pertinent to the research questions were selected as primary studies for the detailed review that follows, while the remaining studies were consulted only as background.

\subsecao{Review of Selected Studies}

The selection process described above yielded the corpus of primary studies analyzed in this review. The following subsections examine in detail the ten selected primary studies, spanning shared-memory multi-core parallelizations of Aho--Corasick (RQ1) and optimizations of the algorithm in security-oriented matching workloads (RQ2); each review summarizes the problem addressed, the proposed technique, the experimental evidence, and its relation to the present design.

\subsubsecao{A Parallel Algorithm of String Matching Based on Message Passing Interface for Multicore Processors}

Qu et al.\ \cite{Qu2016MPI} address the limitations of classical multiple-pattern string matching algorithms, such as the Aho-Corasick algorithm, which were originally designed for single-core systems and often become throughput bottlenecks in modern network intrusion detection systems under heavy traffic. The authors highlight the need to exploit the parallel computation power of modern multicore processors to maintain high performance in deep packet inspection.

To solve this performance bottleneck, the researchers propose a parallel Aho-Corasick string matching algorithm utilizing the Message Passing Interface (MPI) on a shared-memory multicore platform. The core methodology relies on data parallelism: a master thread constructs the automaton sequentially and partitions the input text into contiguous blocks. To prevent missed cross-boundary matches, these blocks overlap by the length of the longest pattern minus one character. The data chunks are then distributed to multiple slave threads via MPI communication, allowing each worker to independently traverse the read-only automaton.

The proposed solution was evaluated on an 8-core multicore system using Snort IDS rulesets of 10,000 and 30,000 patterns. The results demonstrate that the Deterministic Finite Automaton (DFA) variant achieves a peak average throughput of 10.5 Gbps, outperforming the Nondeterministic Finite Automaton (NFA) approach in speed and stability, albeit with higher memory consumption. Furthermore, the evaluation shows that throughput scales monotonically with the addition of CPU cores and thread-level parallelism.

Ultimately, the paper demonstrates that a master-slave data partitioning strategy with overlapping boundaries is highly effective for parallelizing the Aho-Corasick algorithm on multicore processors. However, the reliance on MPI for a single shared-memory node introduces serialized communication overhead, which highlights the potential architectural advantages of utilizing native shared-memory threading models, such as POSIX Threads, for lock-free coordination.

\subsubsecao{A Flexible Pattern-Matching Algorithm for Network Intrusion Detection Systems Using Multi-Core Processors}

Lee and Yang \cite{LeeYang2017FHBM} tackle the throughput bottleneck imposed by Aho-Corasick-based pattern matching in signature-based Network Intrusion Detection Systems (NIDS), where the matching phase can account for up to 70\% of total execution time. The authors observe that, although the canonical Aho-Corasick Deterministic Finite Automaton (AC-DFA) offers linear-time scanning, its state-transition table imposes heavy memory traffic and frequent last-level cache misses as the pattern set grows. The prior Head-Body Matching (HBM) algorithm mitigates this by partitioning the AC-DFA into a fast head (DFA) and a SIMD-accelerated body (NFA), but its rigid depth-based partitioning criterion typically overshoots or undershoots the cache budget, yielding suboptimal throughput.

To overcome this rigidity, the researchers propose the Flexible Head-Body Matching (FHBM) algorithm. Instead of partitioning the automaton at a fixed depth, FHBM employs a greedy procedure that fills the head with up to a configurable maximum number of states (\texttt{HSIZE}), prioritizing parent nodes with the largest number of children and preserving the invariant that all siblings remain on the same side of the partition. The head is represented as a full 256-entry transition table for $O(1)$ lookup, while the body is compacted into an NFA traversed via 128-bit SIMD comparisons. Parallel execution is realized at the thread level: four worker threads, one per physical core, concurrently traverse the read-only Head-Body Finite Automaton (HBFA) on the input stream, exploiting the immutability of the automaton after construction to dispense with synchronization in the search phase.

The evaluation was conducted on a quad-core Intel Core i7-3770 (3.4~GHz, 8~MB L3) running 64-bit Ubuntu 14.04, using pattern dictionaries extracted from the Snort ruleset (8{,}673 patterns) and the ClamAV signature database (5{,}248 and 2{,}899 patterns for the type-1 and type-3 sets, respectively). Synthetic input streams were generated by embedding pattern prefixes into a clean corpus -- the King James Bible (HTML) for Snort and the concatenation of \texttt{/usr/bin} binaries for ClamAV -- with controlled match ratios of 1\%, 8\% and 32\%. Throughput reported in MB/s, averaged over 100 runs, shows that FHBM consistently outperforms HBM: gains of approximately +58\% and +46\% are obtained on Snort and ClamAV type-1 at a 1\% match ratio, and +55\% on ClamAV type-3 at a 32\% match ratio, with peak throughput reaching $\approx 787$~MB/s on Snort and exceeding $1{,}300$~MB/s on ClamAV type-3 when the head transition table still fits within the L3 cache.

Overall, the paper demonstrates that a flexible, state-count-based partitioning of the AC-DFA, combined with SIMD body evaluation and multi-core data parallelism, yields a workable throughput envelope on commodity hardware. It also identifies cache capacity as the dominant performance ceiling: once the active state memory exceeds the L3 budget, additional states no longer improve throughput, a finding that reinforces the cache-locality concerns inherent to any shared-memory parallel implementation of Aho-Corasick. Two architectural choices, however, contrast with the present POSIX Threads approach: the authors do not explicitly describe how the input stream is partitioned across the four worker threads nor how cross-boundary matches are handled, and their speedup is tightly coupled to platform-specific SIMD intrinsics rather than to portable data-parallel chunking with overlap. This positions FHBM as a complementary, automaton-restructuring strategy whose throughput gains stem from microarchitectural exploitation rather than from explicit input partitioning, providing a useful comparative reference point for the present Pthreads-based design.

\subsubsecao{An Optimized Parallel Failure-less Aho-Corasick Algorithm for DNA Sequence Matching}

Thambawita et al.\ \cite{Vajira2018} investigate the throughput limitations of the canonical Aho-Corasick (AC) algorithm in bioinformatics workloads, where the dictionary may contain thousands of DNA patterns and the input may reach hundreds of megabytes per scan. The authors observe that, although the Parallel Failure-less Aho-Corasick (PFAC) library of Lin et al.\ already exposes massive thread-level parallelism by launching one thread per input character, it remains a domain-agnostic implementation whose 256-entry transition table and dispersed memory layout do not exploit two structural properties of DNA matching: the alphabet contains only four symbols ($\{A,T,C,G\}$) and the matching kernel exhibits strong spatial and temporal locality. As a consequence, the baseline PFAC implementation suffers from underused memory bandwidth and poor cache behaviour on the GPU memory hierarchy.

To overcome these limitations, the authors propose an application-specific optimization of PFAC on a NVIDIA Fermi GPGPU comprising three coordinated changes. First, the AC transition table is reduced from a generic $256 \times N$ matrix to a compact $4 \times N$ table aligned with the DNA alphabet, drastically shrinking the working set. Second, the transition table is mapped onto texture memory to exploit its two-dimensional spatial caching, while the input text is staged into shared memory to exploit temporal locality across the threads of each block. Third, the \texttt{nextState} and \texttt{matchedPatternId} arrays are merged into a single contiguous one-dimensional array so that consecutive accesses benefit from cache-line reuse. The underlying parallel decomposition is preserved from PFAC: every input character is assigned its own thread, which independently traverses the failure-less automaton until a mismatch terminates its scan, and matches are recorded in thread-local result slots that require no inter-thread synchronization during the search phase.

The evaluation is conducted on a NVIDIA Tesla C2075 (Fermi architecture), using five synthetic DNA dictionaries (ranging from $1{,}000$ to $5{,}000$ patterns) and five synthetic DNA input datasets sized between 76~MB and 380.2~MB. The reported metric is wall-clock execution time, measured against the unmodified PFAC library; absolute throughput in MB/s and any sequential CPU baseline are notably absent. Across all configurations, the optimized implementation consistently outperforms the original PFAC, with the largest datasets exhibiting approximately a $3\times$ speedup over the baseline library. The authors further observe that memory-layout optimizations -- in particular the merged one-dimensional array and the use of texture memory -- have a substantially larger impact on performance than tuning L1 cache sizes or per-block thread configurations.

This study reinforces two principles that are directly transferable to the present Pthreads-based design even though the target hardware differs: first, that the read-only nature of the AC automaton after construction enables fully lock-free, thread-local matching, irrespective of whether parallelism is realized on a GPU or on a multi-core CPU; second, that the placement of the automaton and the input within the cache hierarchy frequently dominates raw thread-count scaling. Two limitations, however, restrict the direct applicability of these findings to the current work. The PFAC decomposition launches one thread per input byte, which is well suited to GPU SIMT execution but mismatched with the much coarser granularity of multi-core CPUs, where the Pthreads model favours chunk-based data partitioning with overlap regions of $\texttt{max\_pattern\_len}-1$ bytes. Moreover, the alphabet-specific reduction to a $4 \times N$ transition table relies on the small DNA alphabet and cannot be reused for the full 256-symbol byte alphabet of IDS/YARA workloads targeted by the present implementation. The Thambawita et al.\ study thus serves as a complementary point of comparison, framing PFAC-style fine-grained parallelism as the GPU counterpart of the coarse-grained, chunked Pthreads strategy adopted in this thesis.

\subsubsecao{The Diversification and Enhancement of an IDS Scheme for the Cybersecurity Needs of Modern Supply Chains}

Deyannis et al.\ \cite{Deyannis2022Diversification} examine the performance and energy burden imposed by software Aho-Corasick matching when an industrial supply chain node must concurrently execute its primary business workload and inspect real-time network traffic. The authors observe that the canonical Snort IDS, although well-tuned, saturates the CPUs of low-power supply chain assets and therefore introduces a deployment gap: edge nodes either drop packets under load or shift inspection costs upstream, conflicting with the principle of distributed perimeter defence advocated by recent supply chain cybersecurity frameworks. The bottleneck is unambiguously located in the multi-pattern matching kernel, which Snort implements with Aho-Corasick.

To restore line-rate operation without altering the IDS configuration model, the authors offload the AC engine to the FPGA fabric of a low-cost PYNQ~Z1 board (Xilinx Zynq-7000 MPSoC). The fully expanded, failure-free DFA is serialized as a flat two-dimensional integer array indexed by \texttt{dfa[state * 256 + char]}, eliminating pointer-chasing and converting every transition into a single arithmetic lookup -- a layout structurally identical to the read-only \texttt{goto\_tbl} adopted by the present Pthreads implementation. The transition table is then fed into a four-stage hardware pipeline implemented in Vivado HLS and clocked at 410~MHz, while the official Snort software stack continues to run on the co-located ARM Cortex-A9 cores; only the AC matching kernel is replaced. Because the design consumes less than 10\% of the available LUTs, the authors note that two or more copies of the IP module could be instantiated to expose additional spatial parallelism, projecting theoretical $2.8\times$ and $5.6\times$ accelerations that, however, are not empirically validated.

The evaluation employs four publicly available IIoT traffic datasets -- the Malware Traffic Analysis blog capture, DARPA-99, the UNSW ToN\_IoT dataset and the NCC Group honeypot trace -- replayed against the same Snort ruleset, comparing the FPGA-accelerated build against a fully optimized software Snort baseline running on the identical PYNQ board. The reconfigurable design delivers $1.2\times$ to $1.4\times$ acceleration in pure execution time (e.g., $79.4$~s versus $108.9$~s on the DARPA-99 trace) and up to $1.4\times$ higher packet-per-second rates on traces dominated by large packets, while underperforming the software baseline by approximately 10\% on small-packet traces because of interrupt and DMA overhead. The total board power draw of $\approx 13.8$~W contrasts with the $\approx 118$~W of a high-end single-socket server; combined with the up-to-$1.4\times$ acceleration, the authors report an energy consumption up to almost $12\times$ lower than that server. Aggregate throughput in Gbps and the exact size of the loaded ruleset are not disclosed.

In sum, the paper confirms that the matching phase of Aho-Corasick is the dominant cost of signature-based IDS in resource-constrained settings and demonstrates that a hardware-pipelined DFA can recover throughput at a fraction of the energy budget. Three observations transfer directly to the present Pthreads-based design: the use of a fully expanded failure-free DFA stored as a flat $N \times 256$ integer array is equally appropriate for shared-memory CPU traversal, since it preserves $O(1)$ per-character transitions while permitting concurrent read-only access by multiple worker threads; the matching phase rather than the automaton construction dominates the cost, justifying the architectural decision to parallelize only the search loop; and packet-size or input-size sensitivity is a workload effect that any chunk-based parallel implementation should evaluate, since small inputs amortize less of the per-chunk start-up cost. The principal divergence is architectural: while the paper realizes parallelism through hardware pipelining and replicated FPGA IP cores, the present Pthreads design pursues coarse-grained data parallelism through chunked partitioning of the input with overlap regions of $\texttt{max\_pattern\_len}-1$ bytes, exploiting commodity multi-core CPUs and avoiding the reconfigurable-hardware deployment cost.

\subsubsecao{Fast String Searching on PISA}

Jepsen et al.\ \cite{Jepsen2019PISA} target the bandwidth and per-byte compute bottlenecks that arise in disaggregated data centres, where storage nodes carry weak CPUs and every byte of incoming data must traverse PCI-Express to reach a general-purpose processor. The authors observe that pre-filtering large volumes of stored or in-flight data with multi-pattern string matchers is essential to keep downstream analytics tractable, but that an x86 host inspecting every byte does not scale to disk-array or line-rate workloads. Aho-Corasick is identified as the canonical algorithmic foundation for this pre-filter, and the paper poses the question of whether its deterministic finite automaton can be executed entirely within the data plane of a programmable network ASIC, exploiting native pipeline parallelism rather than thread-level parallelism on a CPU.

To answer this question, the authors introduce PPS, a P4-based compiler and runtime that lowers an Aho-Corasick-derived DFA onto the PISA architecture of the Barefoot Tofino switch. Two transformations are applied to the canonical automaton. First, a $k$-stride DFA is synthesized by closing the standard one-character transition function over inputs of $k$ bytes, so that each pipeline pass consumes several characters at the cost of a multiplicatively larger transition table; this trades on-chip SRAM against per-stage throughput. Second, the resulting DFA is optionally partitioned into up to three independent sub-automata that are placed in different pipeline stages and traversed in parallel against the same packet, allowing the aggregate pattern set to exceed the memory of any single stage. To enforce correctness when packets are processed independently, the system mandates that adjacent input chunks overlap by the length of the longest pattern -- a constraint conceptually identical to the $\texttt{max\_pattern\_len}-1$ overlap adopted by the present Pthreads implementation, although applied here at the network MTU granularity. An optional CRC16 hash representation of the transition table is offered as an accuracy-for-memory trade-off; this approximation is incompatible with exact-match requirements and is therefore restricted to use cases tolerant of false positives.

The evaluation contrasts PPS running on a Tofino switch against two CPU baselines -- Spark SQL on a four-node dual-socket Intel Xeon E5-2603 cluster and Unix \texttt{grep} on a 12~GB RAM-disk log -- using Snort community-rule \texttt{content} patterns of up to 32 characters and three corpora: the 4.7~GB Podesta e-mail archive, a 212~GB Twitter streaming capture, and the synthetic log file. End-to-end execution times measured at 128 patterns yield a $>20\times$ speedup over Spark SQL on the e-mail workload and a $6.5\times$ speedup on the tweet workload (5.43~min versus 35.4~min), while throughput at 100 Snort patterns with stride~4 is calculated -- not measured -- at $3.8$~Tbps for a 64-port Tofino. The authors are explicit that the Tbps figure is a theoretical upper bound derived from the deterministic line rate of the compiled P4 program; CPU baselines are measured end-to-end, and external GPU/FPGA references (e.g., 150~Gbps on GPU, 45~Gbps for DFC on x86) are quoted rather than reproduced.

The paper substantiates Aho-Corasick as the algorithmic foundation that scales from a single CPU core all the way to multi-Tbps switching silicon, and it makes two design decisions that map directly onto the present Pthreads-based work. The mandatory overlap of \emph{longest-pattern} bytes between adjacent input chunks is the same correctness primitive that the present implementation enforces by reserving $\texttt{max\_pattern\_len}-1$ warm-up bytes at every chunk boundary; this independent confirmation is particularly valuable because it derives from an entirely different parallel substrate, ruling out the possibility that the overlap is an artifact of thread-level decomposition. Likewise, DFA partitioning provides a precedent for splitting a pattern set across independent execution units, although the present work partitions the \emph{input} rather than the automaton, leaving the goto table intact for shared, read-only access by all Pthreads workers. The principal divergence is the hardware target: PPS realizes parallelism through replicated pipeline stages and TCAM/SRAM primitives in a switch ASIC, whereas the present design pursues coarse-grained data parallelism on commodity multi-core CPUs, accepting lower theoretical throughput in exchange for portability and exact matching.

\subsubsecao{Practical Multi-Pattern Matching Approach for Fast and Scalable Log Abstraction}

Tovar{\v{n}}ák \cite{Tovarnak2016REtrie} confronts the throughput limitations of multi-pattern matching in log-abstraction pipelines, where each incoming log line must be classified by matching it against a potentially large set of regex-like patterns. The author observes that the naive strategy of iterating over the entire pattern set per log line incurs $O(|\text{patterns}|)$ cost per input and degrades severely once the dictionary contains more than a few tens of entries, while general-purpose multi-regex libraries such as Google RE2 hit memory-cap limits when the DFA representation exceeds a few hundred patterns. The work is positioned alongside Aho-Corasick as a multi-pattern matcher built on a trie-based automaton, but it explicitly targets the same parallelism model adopted by the present thesis: a single automaton constructed once and shared, read-only, across multiple concurrent workers on a commodity multi-core CPU.

To recover throughput, the author proposes REtrie, a compressed trie that combines literal-character pointers with regex \emph{token} pointers (named wildcards) and supports controlled backtracking only over the token segments, preserving an $O(|\text{input}|)$ matching cost in the typical case. Three design choices map directly onto the present Pthreads design: the trie is constructed sequentially and then declared immutable, so that worker processes need no synchronization during the matching phase; node arrays are compressed with PATRICIA-style path compression and two-level adaptive nesting, reducing the memory footprint of the otherwise sparse 256-entry pointer arrays and improving cache locality; and data-level parallelism is realized by routing disjoint subsets of input log lines to independent Erlang OTP processes, each of which traverses the same shared automaton without any inter-process locking. The system is implemented in Erlang with native-code generation via the HiPE compiler, so the parallelism model is process-based rather than thread-based, but the architectural invariant -- read-only shared automaton, thread-local results, no locks on the hot path -- is the same one adopted by the Pthreads design of the present thesis.

The evaluation is conducted on a dual-socket Intel Xeon E5-2650~v2 (2.60~GHz) with 64~GB of RAM, scaled from one to eight cores. The pattern sets include a hand-curated set of 22 syslog \texttt{auth/authpriv} patterns ($P_{22}$) and large Apache Hadoop sets ($R_{x}$) with up to 3,500 entries derived from source-code analysis. The runtime workloads include $D_0$, a real-world capture of 100 million syslog lines collected over one month, plus partially synthetic derivatives ($D_x$, $H_x$) of 1.1 million messages each constructed by controlled duplication and shuffling. Each measurement is averaged over ten runs with extreme values discarded. The single-core comparison shows that REtrie's running time remains essentially constant as the pattern count grows from 1 to 22, whereas the naive baseline reaches roughly 12~s on $P_{22}$ versus ${\approx}3$~s for REtrie and RE2; for $R_x \geq 2{,}000$ patterns the Erlang REtrie even surpasses C++~RE2, which collides with its default DFA memory limit. The multi-core scaling on $D_0/P_{22}$ exhibits a near-ideal $2\times$ speedup at two cores; at eight cores the 100-million-line workload is processed in $52.543$~s, equivalent to ${\approx}1{,}903{,}203$ log lines per second, with the speedup growth flattening past four cores, in line with the ceiling expected from Amdahl's law \cite{Amdahl1967} and from memory bandwidth.

Taken together, the paper provides independent empirical validation of two design choices that underpin the present Pthreads-based implementation of Aho-Corasick. First, the architectural invariant of a read-only shared automaton accessed concurrently by multiple worker processes is shown to deliver near-ideal scaling at low core counts on commodity multi-core hardware, with the diminishing-returns regime located around four to five cores, an observation consistent with the early plateau found in the present experiments. Second, the use of path compression and adaptive node sizing to improve cache locality reinforces the importance of the flat \texttt{goto\_tbl} layout employed by the present implementation, even though the algorithmic context differs. Two limitations restrict the direct transferability of the results. The REtrie automaton is not an Aho-Corasick DFA: it lacks the failure and dictionary-suffix functions and processes each log line as a single bounded record, so the chunk-boundary overlap problem that motivates the present implementation's $\texttt{max\_pattern\_len}-1$ warm-up region simply does not arise here. Moreover, Erlang OTP scheduling abstracts away the cache, false-sharing and NUMA effects that the Pthreads design must address explicitly. The Tovar{\v{n}}ák study therefore serves as a precedent for the parallel \emph{strategy} (lock-free shared immutable automaton) rather than for its concrete realization, and as a useful throughput reference point on comparable Xeon-class hardware for multi-pattern matching at the line-of-text granularity.

\subsubsecao{Pattern Matching in YARA: Improved Aho-Corasick Algorithm}

Reg{\'e}ciov{\'a} et al.\ \cite{Regeciova2021MYARA} address a particular failure mode of the canonical Aho-Corasick implementation embedded in the YARA rule-matching framework, a representative production scanning tool that relies on this algorithm. The authors observe that YARA's preprocessing phase extracts short literal substrings -- termed \emph{atoms} -- to seed the AC automaton, but the atom-selection heuristic fails to represent the wildcards, character classes and metacharacters that increasingly appear in industrial signatures. As a result, the scanner either selects empty atoms, in which case the matching loop degenerates into a naive byte-by-byte comparison, or it issues a ``slowing down'' warning that renders the rule unusable on large-scale platforms such as VirusTotal Hunting. The bottleneck is therefore located not in raw matching throughput but in the construction-time choice that determines which literal fragments seed the AC automaton.

To overcome this limitation, the authors propose MYARA, a fork of YARA~4.0.2 introducing three coordinated changes to the preprocessing phase while leaving the matching loop intact and sequential. First, atoms are encoded as 256-bit \emph{bitmaps}, where each bit indicates whether the corresponding byte value is admissible at that position; a literal character, the dot metacharacter and a class such as \texttt{[a-z]} are therefore expressed in a uniform structure that supports fast subset and inclusion tests through bitwise operators. Second, the AC \texttt{goto} function is generalized to accept \emph{sets} of characters per transition rather than single bytes, and the resulting non-determinism is resolved by an explicit determinization step before the failure function is computed. Third, the existing Aho-Corasick Interleaved State Matrix (ACISM) -- a flat array of 32-bit integers that encodes both \texttt{goto} transitions and failure links for $O(1)$ lookup -- is updated to accommodate the denser slot occupancy produced by multi-symbol transitions. Because the resulting automaton remains entirely read-only after construction, the same shared-memory invariant relied upon by the present Pthreads design continues to hold.

The evaluation is conducted on a server equipped with an Intel Xeon E5-2697~v4 at 2.30~GHz running Debian 4.9.168, with binaries built using GCC~8.3.0. Three pattern dictionaries of contrasting composition are exercised: a hand-crafted set of five YARA rules targeting Bitcoin addresses, e-mail addresses and the INI format; the publicly available ReversingLabs corpus comprising 371 rules with $0\%$ regular expressions; and the proprietary Avast set of 13{,}026 rules, of which only $2.1\%$ are regular expressions. All workloads are scanned against an 8.54~GB real-world corpus that aggregates Bitcoin transaction history, e-mail archives and INI configuration files; each reported measurement is the minimum of fifteen daily runs collected over six days. On the degenerate Bitcoin regex (v1) the scan time falls from $169.700$~s under canonical YARA to $17.623$~s under MYARA, i.e.\ a ${\approx}9.6\times$ speed-up that recovers the rule from the unusable regime. On the realistic ReversingLabs ruleset, where atoms are dominated by literal and hexadecimal strings, scanning time is essentially unchanged ($1{,}503.65$~s vs $1{,}504.90$~s); on Avast the global improvement is approximately $1\%$. Throughput in MB/s and memory consumption are not reported.

On balance, the paper offers two contributions that bear directly on the present Pthreads-based work even though it introduces no parallelism of its own. First, it documents and characterizes the ACISM representation that the present implementation -- and any other shared-memory parallel Aho-Corasick variant -- can adopt as a strictly read-only substrate, with the additional benefit that bitmap-encoded transitions remain compatible with the flat \texttt{goto\_tbl[state * 256 + byte]} layout adopted in this thesis. Second, it provides documented single-threaded baselines on a 8.54~GB realistic corpus that serve as external reference points for the present evaluation. Three limitations restrict the direct applicability of the results, however: the optimization is concentrated in the construction phase, whereas the thesis specifically targets the matching phase; the reported gains evaporate on rulesets dominated by literal patterns (the ReversingLabs case), which closely resemble the IDS and YARA workloads exercised in the present evaluation; and the absence of throughput measurements in MB/s precludes a direct quantitative comparison. The Reg{\'e}ciov{\'a} et al.\ study therefore complements rather than competes with the present work: it improves the automaton fed to the scanner, while the present thesis improves the parallel execution of the scanner against an unchanged automaton.

\subsubsecao{SIMD-Accelerated Regular Expression Matching}

Sitaridi, Polychroniou and Ross \cite{Sitaridi2016SIMD} focus on the sequential bottleneck that arises when a DFA-based pattern matcher is embedded in an in-memory database column-filter operator: the per-character transition $s_{t+1} = \delta(s_t, c_t)$ creates a tight data dependency that defeats straightforward SIMD auto-vectorization, leaving most lanes of a wide vector register idle. The authors explicitly observe that Aho-Corasick is equivalent to a minimal DFA once the failure transitions are absorbed into the transition table, so the analysis of vectorized DFA traversal performance transfers directly to the matching phase of any AC engine. Their concern is therefore the same matching-phase throughput targeted by the present thesis, but addressed at a different level of the parallelism hierarchy.

To exploit the unused lanes, the paper introduces a non-lockstep, short-circuit SIMD-DFA traversal algorithm. Rather than advancing all $W$ lanes of the SIMD register in unison one character at a time -- which wastes work whenever a string is rejected or accepted early -- the algorithm tracks an independent input offset and DFA state per lane, uses SIMD \emph{gather} instructions to fetch the relevant byte for each lane in a single instruction, and replaces any lane that reaches a sink state with a fresh input through selective loads. Three auxiliary techniques sustain the throughput: an eight-byte buffered gather amortizes the cost of non-contiguous loads, a two-way unrolled inner loop hides instruction latency, and a branch-free selective store records the identifiers of accepted strings. Critically for the present discussion, the authors compose this SIMD-level data parallelism with thread-level parallelism by spawning one worker per hardware thread; the DFA transition table is shared, read-only, across all workers, and string offsets are partitioned among them. This is structurally identical to the present Pthreads design, with the difference that the thesis partitions a \emph{single} large input into chunks, whereas the paper partitions a column of many short strings.

The evaluation contrasts two platforms: an Intel Xeon E3-1275v3 with four Haswell cores and 256-bit AVX2, and an Intel Xeon Phi 7120P with 61 cores and 512-bit SIMD. Pattern dictionaries include a 90-state URL-validation regex (23~KB transition table), a 9-state e-mail regex, and synthetic multi-pattern dictionaries built from random English words whose DFA size is swept from $10$ to $10^{5}$ states to cross the L1, L2 and last-level cache boundaries. Inputs are synthetically generated string columns of varying length, selectivity and average failure point. Throughput is reported in GB/s of total bytes scanned together with the fraction of peak DRAM bandwidth attained. The non-lockstep algorithm yields up to $2\times$ over a scalar baseline on the Haswell CPU and up to $5\times$ on the Xeon Phi, with the gap over lockstep reaching $3\times$ for long strings that reject early. Two further observations are particularly relevant. First, throughput scales linearly with hardware-thread count on both platforms, confirming that SIMD-level and thread-level parallelism compose without contention as long as the DFA is shared read-only. Second, performance peaks when the DFA fits in the L1/L2 cache: once the transition table grows past last-level cache, vectorization gains shrink, with throughput reverting to a memory-bandwidth-bound regime.

Ultimately, three findings carry directly into the present Pthreads-based work, with the caveat that the underlying parallelism mechanisms differ. First, the explicit observation that AC is equivalent to a minimal DFA, combined with the demonstrated linear scaling of thread-level parallelism over a shared read-only transition table, independently validates the architectural choice of immutable post-construction \texttt{goto\_tbl} central to the present design. Second, the empirical demonstration that cache residency dominates throughput once the DFA exceeds last-level cache provides a documented mechanism behind the saturation that the present thesis observes on large rule sets, motivating both the cache-locality concerns of the flat \texttt{goto\_tbl[state * 256 + byte]} layout and the evaluation of memory-bandwidth contention as the thread count grows. Third, the orthogonality of SIMD-level and thread-level parallelism documented in this paper makes it a natural reference for positioning the Pthreads strategy as a baseline that future work could combine with SIMD intrinsics in the inner loop. Two limitations bound the transferability of the numerical results: the workload is restricted to fixed-length strings in database columns, in contrast to the variable-length chunked file scans targeted by the present thesis, and the cache-residency assumption holds only for the small URL regex used in most experiments and breaks down for the large multi-pattern dictionaries that motivate AC in security applications.

\subsubsecao{Harry: A Scalable SIMD-based Multi-literal Pattern Matching Engine for Deep Packet Inspection}

Xu et al.\ \cite{Xu2023Harry} tackle the per-byte cost of multi-literal pattern matching in software Deep Packet Inspection (DPI), which the authors identify as the dominant bottleneck of production stacks such as Snort, Suricata, ModSecurity and CloudFlare's WAF. The canonical Aho-Corasick (AC) algorithm is acknowledged as the traditional choice but is shown to leave large performance margins on the table compared with the SIMD-accelerated FDR engine of Hyperscan; FDR itself, however, requires three SIMD operations per input character regardless of the vector width, attains only $50\%$ SIMD-lane utilization, and depends on the \texttt{VPSLLDQ} shift instruction whose 128-bit lane boundary makes AVX-512 unusable. The paper therefore poses the question of how far the throughput ceiling of single-threaded multi-literal matching can be raised on commodity x86 CPUs by aggressively redesigning the SIMD kernel, establishing a useful upper bound against which the multi-core Pthreads strategy of the present thesis can be positioned.

Three coordinated techniques constitute the Harry engine. A column-vector-based shift-or matching algorithm replaces FDR's row-vector decomposition: instead of loading one row per input character, Harry loads one column vector per literal-position slot, so that the number of SIMD operations becomes proportional to the literal length $\ell$ (fixed at eight in the implementation) rather than to the input length, yielding only $0.41$--$0.71$ operations per input character on AVX-512 and raising vector utilization to $87.5\%$. Two mask-table encodings -- a compression-based variant (\textit{Harry6b}, optimal for small rule sets) and a decomposition-based variant (\textit{Harry12b}, optimal for large rule sets) -- are selected at compile time by a heuristic based on the zero-bit rate of each encoding, keeping the column vector within 512 bits in both regimes. Finally, the legacy \texttt{VPSLLDQ} shift is replaced with a \texttt{VPERMB}-based shuffle, eliminating cross-lane data loss and enabling the kernel to exploit AVX-512 without algorithmic regressions. The matching pipeline preserves Hyperscan's two-stage structure -- approximate filter followed by exact verification -- so the read-only nature of the underlying tables remains intact, consistent with the lock-free shared-automaton invariant adopted by the present Pthreads design.

Harry is evaluated on an Intel Xeon Gold 6348 (2.60~GHz, two sockets of 28 cores each, AVX-512) running Ubuntu 20.04, with GCC~9.4.0 used without optimization flags. Literal rules are extracted from two production rule sets -- OWASP ModSecurity Core Rule Set v3.3.2 and Snort Emerging Threats v2.9.0 -- and partitioned into nine subsets spanning $10$ to $3{,}000$ rules. Three input traces are exercised: $121$~MB of real IXIA HTTP packets, $4$~MB of Alexa non-HTTP packets and $763$~KB of synthetic random bytes; each experiment is averaged over $1{,}000$ runs. The reported throughput ranges from $30$ to $70$~Gbit/s and represents a $26\times$--$45\times$ speed-up over canonical Aho-Corasick on real HTTP and non-HTTP traffic, peaking at $52\times$ on the random byte stream, and a $1.06\times$--$2.09\times$ improvement over FDR. The heuristic encoding selector picks the better variant in every evaluated configuration, and Harry maintains its advantage across the full $10$--$3{,}000$ rule range, in contrast to the prior Teddy engine, which degrades sharply beyond roughly sixty rules.

Overall, the paper has two effects on the related-work positioning of the present thesis. First, it concretely quantifies the ceiling that single-threaded SIMD literal matching can reach on commodity AVX-512 hardware -- $30$--$70$~Gbit/s at $0.41$ operations per character with $87.5\%$ vector utilization -- and reports that this regime is up to $52\times$ faster than the canonical Aho-Corasick algorithm. This establishes Harry as the contrast point against which a coarse-grained Pthreads parallelization of AC must justify itself, and clarifies that the rationale of the present work is not to outperform Harry on a single thread but to exploit multi-core resources where the SIMD path is already saturated, since the same input can be partitioned across cores and the AC automaton kept fully shared and read-only. Second, the paper documents that SIMD-level and thread-level parallelism are architecturally orthogonal: Harry runs entirely within a single thread, so its gains compose multiplicatively with any data-parallel chunking strategy on multi-core CPUs. Two caveats restrict the direct numerical comparability of the results: the AC baseline is compiled by GCC without optimization flags while the Harry path is hand-vectorized, which likely inflates the reported $52\times$ ratio, and the engine replaces rather than parallelizes the Aho-Corasick automaton, departing from the strict-AC requirement of the present thesis. Harry should therefore be cited as the canonical ceiling for single-threaded literal matching in DPI, not as a parallel AC variant.

\subsubsecao{Characterizing Realistic Signature-based Intrusion Detection Benchmarks}

Aldwairi, Alshboul and Seyam \cite{Aldwairi2018Characterizing} address a methodological gap that affects every empirical study of pattern matching for signature-based Intrusion Detection Systems (NIDS), including the present thesis: there is no agreed standard for the datasets, the metrics or the simulation environment under which Aho-Corasick and competing algorithms are evaluated, and the heterogeneity that results from this absence renders cross-paper comparisons unreliable. The authors review the most cited public datasets -- DARPA, KDD~Cup~99, NSL-KDD, Kyoto~2006+ and CAIDA -- and document that all of them were designed for anomaly-based detection and therefore lack representative attack signatures; privately collected traces, conversely, cannot be redistributed for legal reasons. The paper therefore positions itself not as a parallelization contribution but as the construction of a reproducible NIDS benchmark stack that subsequent works -- this thesis included -- can adopt to ground throughput and speed-up claims in realistic conditions.

The contribution is delivered as a five-part artifact. First, a GUI-based parser ingests heterogeneous IDS rule formats and extracts the \texttt{content} and \texttt{uricontent} keywords of 120 Snort rule files, yielding hexadecimal-encoded signature dictionaries of varying cardinality (from 300 to over 2{,}400 signatures, suitable for scalability sweeps). Second, eight representative traces are curated from the 78 DEFCON17 Capture-the-Flag captures: four ``good'' traces (audio-video streaming, downloads, Hotmail, live chat) representing baseline traffic, two ``ugly'' traces (51 and 58) where 45\% and 43\% of packets carry attack signatures, and two ``bad'' traces (1 and 22) with intermediate 16\%--17\% malicious densities. Third, a pluggable matching engine implemented as a Windows DLL allows researchers to substitute their own algorithm under identical input and instrumentation. Fourth, four classical algorithms -- Aho-Corasick, Knuth-Morris-Pratt, Wu-Manber and Rabin-Karp -- are reimplemented strictly sequentially against the same engine to provide single-thread reference numbers. Fifth, an explicit set of metrics is fixed: runtime in microseconds, memory in bytes, and count of matched attack signatures per trace, each averaged over ten repetitions.

The hardware platform is not disclosed; the engine is implemented in C\# and the authors flag that the runtime's garbage collector limits the precision of memory measurements. The reported numbers are nevertheless useful as a sequential baseline. Aho-Corasick is the slowest of the three viable algorithms on the large ``audio-video'' trace, taking $2{,}374{,}992$~µs against Wu-Manber's $619{,}469$~µs, and consumes the largest memory footprint, peaking at $6{,}840{,}694$~bytes versus $2{,}694{,}016$~bytes for KMP. The runtime-versus-signature-count sweep on the same trace shows that AC scales non-linearly: the time grows from $168{,}079$~µs at $300$ signatures to $2{,}374{,}992$~µs at the full dictionary, an effect the authors describe as ``almost exponential''. The memory-versus-packet-count sweep on the ``ugly'' trace~58 reports a similar growth -- from $29{,}402$~bytes at $25{,}000$ packets to $3{,}871{,}408$~bytes for the full trace -- although this expansion most likely reflects accumulated match-record storage and C\# allocation behaviour rather than an intrinsic growth of the AC automaton, whose DFA trie is fixed by the signature set.

In sum, the paper provides the present thesis with two contributions that no other reviewed work delivers in combination. First, the DEFCON17 trace selection and the Snort-derived signature dictionaries are documented, redistributable benchmark inputs that can be replayed against the present Pthreads implementation to expose the multi-core speed-up under the same workload conditions used elsewhere in the literature. Second, the per-trace sequential AC runtimes -- in the order of tens of milliseconds for the smaller traces and a few seconds for the largest -- furnish numerical reference points against which the present multi-thread throughput can be situated. Three limitations restrict the direct use of the published numbers. The hardware on which they were measured is not specified, so absolute runtimes cannot be reproduced. The implementation language (C\#) and its garbage collector make memory figures imprecise. And the paper does not report throughput in MB/s or Gbps, which are the metrics the present thesis adopts; conversion from microseconds plus payload size in bytes is therefore required when quoting the Aldwairi figures as competitors. The paper should be cited for benchmark methodology and realistic-workload selection rather than as a parallel-AC technique, and it complements the multi-core, multi-threaded contributions reviewed above by anchoring the evaluation in the canonical NIDS context that motivates Aho-Corasick parallelization in the first place.

% Fim do conteúdo da seção de Revisão Sistemática.